\section{Completion and the stochastic block model}\label{sec:completion}
Table~\ref{tab:coverage} collects the arguments for the equality cases.
In the table, $a=|\lambda|$ denotes the negative channel magnitude.

\begin{table}[ht]
\centering\small
\begin{tabular}{@{}p{.21\textwidth}p{.36\textwidth}p{.34\textwidth}@{}}
\toprule
Equality case & Parameter range & Argument\\
\midrule
Positive channel & All admissible $d>1$ & Analytic exclusion, Section~\ref{sec:positive}\\
Negative, regular & $b\ge5$, so $a=b^{-1/2}<.448$ & Weighted identities, Section~\ref{sec:weighted}\\
Negative, regular & $b=4$, $a=1/2$ & Finite domination, Section~\ref{grid:section}\\
Negative, Poisson & $a=d^{-1/2}\le.477$ & Weighted identities, Section~\ref{sec:weighted}\\
Negative, Poisson & $4\le d\le881/200$ & Finite domination, Section~\ref{grid:section}\\
\bottomrule
\end{tabular}
\caption{Coverage of $d\lambda^2=1$, with $a=|\lambda|$ on the
negative side. The two Poisson ranges overlap since
$881/200>(1000/477)^2$. Below branching parameter four, negative channels
follow by degradation from the coloring case.}\label{tab:coverage}
\end{table}

\begin{proof}[Proof of Theorem~\ref{thm:threshold}]
Proposition~\ref{fnd:supercritical} proves reconstruction whenever
$d\lambda^2>1$.  Conversely, Proposition~\ref{pos:nonreconstruction}
proves non-reconstruction for $0\le\lambda\le d^{-1/2}$, and
Proposition~\ref{grid:negative-completion} proves it for every admissible
negative parameter with $d\lambda^2\le1$.  These conclusions include
all equality cases and exhaust the channel domain.  The Poisson
statements use the original tree law, without conditioning on survival.
\end{proof}

We return to the stochastic block model of Section~\ref{sec:sbm},
where $a,b>0$ are the within- and between-community connectivity
parameters and $d=(a+2b)/3>1$.  The remaining step is to apply the
external tree-to-graph and algorithmic results to the tree theorem.

\begin{proof}[Proof of Corollary~\ref{cor:sbm}]
Set
\[
 \lambda=\frac{a-b}{a+2b}\in(-1/2,1),\qquad
 d\lambda^2=\frac{(a-b)^2}{3(a+2b)}.
\]
The corresponding Poisson broadcast process has exactly the transition
matrix~\eqref{eq:channel}. Suppose first that $d\lambda^2\le1$.
Theorem~\ref{thm:threshold} gives $a_\ell\to0$ for this process.

To match the fixed-tree formulation of non-reconstruction, let $T$ be
the entire Poisson tree and put
\[
 A_\ell(T)=\E\!\left[
 \left\|\PP(\sigma_\rho=\cdot\mid T,\sigma_{\partial T_\ell})
       -\operatorname{Unif}\{1,2,3\}\right\|_{\mathrm{TV}}
 \mathrel{\Big|}T\right].
\]
For almost every fixed $T$, the boundary experiment is degraded as $\ell$
increases, so $A_\ell(T)$ decreases to a nonnegative limit $A_\infty(T)$.
The offspring beyond depth $\ell$ are independent of the root and
depth-$\ell$ spins conditional on $T_{\le\ell}$, and hence
$\E A_\ell(T)=a_\ell$. Bounded convergence therefore gives
$\E A_\infty(T)=0$, so $A_\infty(T)=0$ for almost every $T$.
If $\mu_{\ell,i}^T$ is the boundary law conditional on root state $i$
and $\overline\mu_\ell^T=\frac13\sum_i\mu_{\ell,i}^T$, Bayes' formula gives
\[
 A_\ell(T)=\frac13\sum_{i=1}^3
   d_{\mathrm{TV}}(\mu_{\ell,i}^T,\overline\mu_\ell^T),
\]
where $d_{\mathrm{TV}}$ denotes total-variation distance. The triangle inequality now shows that the pairwise total-variation
distances between the $\mu_{\ell,i}^T$ tend to zero for almost every $T$.
This is the non-reconstruction hypothesis of the tree-to-graph transfer
of Mossel, Sly, and Sohn \cite[Theorem~2.1]{MosselSlySohn2025}.
Their conclusion, with the overlap of their Definition~1.4, gives
\[
 \limsup_{n\to\infty}\E\operatorname{ov}(\sigma,\widehat\sigma_n)\le\frac13
\]
for deterministic estimators. Any violation by randomized estimators would,
for each $n$, admit a deterministic choice of the independent seed with
expected overlap at least the randomized expectation minus $1/n$,
contradicting the same bound. For every pair of labelings, the average
agreement over all six permutations is $1/3$, so its maximum is at least
$1/3$. This proves (i).

If $d\lambda^2>1$, Corollary~1 of Abbe and Sandon
\cite{AbbeSandon2018}, with three communities, gives an $O(n\log n)$
algorithm for detection. Their Lemma~1 converts their two-set detection
criterion, for uniform community probabilities, into overlap exceeding
$1/3$ by a fixed positive amount with probability tending to one.
This proves (ii).
\end{proof}

The corollary covers both $a>b$ and $a<b$. Its positive-parameter
hypothesis is part of the statement: the boundary cases $a=0$ or $b=0$
are not included. In particular, the negative equality has $d>4$ here;
the coloring equality $a=0$, $b=6$, $d=4$ is excluded.

